\section{One retained weight for the matched cubic family at coefficient seven}
\label{sec:mixed-cubic-residual-trivial-higher-lifting}

Put $K=\mathbb Q(\omega)$, $t^3=269$, and $L=K(t)$.
Let $\chi$ be either faithful character of $\operatorname{Gal}(L/K)$.
Its values lie in $\mu_3\subset\mathbb Z_7^\times$.
Write
\[
 c_7=\omega_7\langle c_7\rangle,\qquad
 \omega_7:G_K\longrightarrow\mu_6\subset\mathbb Z_7^\times,
 \qquad \langle c_7\rangle\in1+7\mathbb Z_7.
 \tag{MRT.1}
\]
Here $\omega_7$ is the Teichm\"uller lift of the cyclotomic reduction.
Let $O$ be the original ring of integers of any finite extension
$E/\mathbb Q_7$, and let $b=\operatorname{rank}_{\mathbb Z_7}O$.
In particular, $O$ may be ramified.  Choose the actual faithful
conductor-seven character $\xi_7$ with values in $\mu_6$ in this $O$.
For $j\in\{1,2\}$, $a\in\{0,\ldots,5\}$, and an integer $e\geq1$,
consider the original coefficient
\[
 D_{j,a,e}=(E/O)(\chi^j\xi_7^a c_7^{-e}),\qquad
 \rho_{a,e}=\overline{\xi}_7^{\,a}\overline c_7^{-e}.
 \tag{MRT.2}
\]
All rational two-power and auxiliary $q,d$ factors are trivial.
The strict group $\mathcal S_7(K,D)$ is the subgroup of
$H^1(G_K,D)$ with zero restriction to the full decomposition group
at both places over seven and to inertia at every other finite place.
Use the same local conditions for finite coefficients.

\begin{theorem}[AC-G219: original-integral matched cubic cost and one retained weight]
\label{AC-G219}
There is an intrinsic set
\[
             \mathscr E\subset\{1,2,3,4,5,6\},
             \qquad |\mathscr E|\leq1,                         \tag{MRT.3}
\]
independent of $O$, the choice of $\chi$ or $\xi_7$, and the host
parameters, with the following properties.  For every coefficient
in (MRT.2) satisfying the literal equality $\rho_{a,e}=1$,
\[
                  \mathcal S_7(K,D_{j,a,e})[7]\simeq O/7O.
 \tag{MRT.4}
\]
If also $e\bmod7\notin\mathscr E$, then the entire strict group is
\[
       \mathcal S_7(K,D_{j,a,e})\simeq O/7O,
       \qquad \frac{\log\#\mathcal S_7(K,D_{j,a,e})}{b}=\log7.
 \tag{MRT.5}
\]
These are assertions on the original coefficient and its full local
conditions.  They hold identically for the actual groups (FC.2).
In particular, every matched weight divisible by seven satisfies
(MRT.5).  The set $\mathscr E$ is defined by the fixed cochain
obstruction in Step 5; its membership or emptiness is not computed
here.  The all-level conclusion below instead retains only one possibly
large weight in each finite cutoff.

For an integer $T\geq2$, put $N=T-1$ and define the actual family
\[
 \begin{split}
 \mathcal A(T)=\{(7,\chi^j\xi_7^a,e):\ &j=1,2,
       \ 0\leq a\leq5,\ 1\leq e<T,\\
       &\rho_{a,e}=1,\quad e\bmod7\notin\mathscr E\},\\
 n_{\mathscr E}(N)&=\sum_{r\in\mathscr E}
             \max\left(0,1+\left\lfloor\frac{N-r}{7}\right\rfloor\right),\\
 |\mathcal A(T)|&=2\bigl(N-n_{\mathscr E}(N)\bigr).
 \end{split}                                                   \tag{MRT.6}
\]
In particular $n_{\mathscr E}(N)\leq\lceil N/7\rceil$.
For any retained set $J$ of original triples in (CP.3), put
$J'=J\setminus\mathcal A(T)$.  The exact replacement of the
nonnegative arithmetic remainder is
\[
 \mathcal R_J=\mathcal R_{J'}+\Gamma_J,
 \qquad \Gamma_J=|J\cap\mathcal A(T)|\log7,
 \qquad 0\leq\Gamma_J<6\mathcal P(F,T).                         \tag{MRT.7}
\]
Here $\mathcal P$ is exactly (CP.2).  Thus a comparison
$|\mathcal I_J-\mathcal R_J|\leq E_{\rm old}+C'\mathcal P$
gives the pointwise bound
\[
 |\mathcal I_J-\mathcal R_{J'}|
       \leq E_{\rm old}+(C'+6)\mathcal P.                       \tag{MRT.8}
\]
The same statement holds with the original strict sum in place
of $\mathcal I_J$.  The already established exponent
$806+16\varepsilon$ for the polynomial envelope is unchanged;
the multiplicative comparison coefficient is enlarged.
The strict groups in (MRT.5) are nonzero; the individual derived
indices retain their separate comparison.
There is no additional window choice or repeated finite-gluing sum.

There is also a completely specified formulation which requires no
knowledge of $\mathscr E$.  For $r=1,\ldots,6$ let
\[
 \begin{split}
 \mathcal A_r(T)=\{(7,\chi^j\xi_7^a,e):\ &j=1,2,\ 0\leq a\leq5,
       \ 1\leq e<T,\\
       &\rho_{a,e}=1,\ e\not\equiv r\pmod7\},\\
 \mathcal R^{\max}_J&=\max_{1\leq r\leq6}
                       \mathcal R_{J\setminus\mathcal A_r(T)}.
 \end{split}                                                   \tag{MRT.24}
\]
These are six fixed, explicitly defined remainders on the same host.
They satisfy
\[
 0\leq\mathcal R_J-\mathcal R^{\max}_J
       \leq2(T-1)\log7<6\mathcal P(F,T).                       \tag{MRT.25}
\]
Consequently (MRT.8) holds with $\mathcal R^{\max}_J$ in place
of $\mathcal R_{J'}$.  A bound on this single maximum is a
sufficient arithmetic gate with the displayed added cost; it
does not ask for six separately existential windows.

The same actual lifting problem has a stronger conclusion at every
power of seven.  There are positive integers $s(e)$, $e\geq1$,
independent of $j$, the original coefficient ring and the host, such
that for every matched label
\[
 \mathcal S_7(K,D_{j,a,e})\simeq O/7^{s(e)}O,
 \qquad b^{-1}\log\#\mathcal S_7(K,D_{j,a,e})=s(e)\log7.
 \tag{MRT.27}
\]
For every $n\geq2$, the set of weights with $s(e)\geq n$ is
empty or one residue class $E_n$ modulo $7^{n-1}$; these residue
classes are nested and $E_2=\mathscr E$.
For $N\geq1$, choose a weight $e_*$ maximizing $s(e)$ on
$\{1,\ldots,N\}$.  Then
\[
 \begin{split}
 \sum_{\substack{1\leq e\leq N\\ e\ne e_*}}s(e)
 &\leq(N-1)+\sum_{k\geq1}\left\lfloor\frac{N-1}{7^k}\right\rfloor\\
 &=(N-1)+v_7((N-1)!)\leq\frac76(N-1).
 \end{split}
 \tag{MRT.28}
\]
No bound on $s(e_*)$ is asserted.

Here is a completely specified arithmetic remainder using this
stronger bound.  Let $\mathcal T(T)$ contain both matched original
labels at every weight $1\leq e<T$, and for $1\leq r\leq N$ put
\[
 \begin{split}
 \mathcal T^{\ne r}(T)&=\{(7,\chi^j\xi_7^a,e):\ j=1,2,
       \ 0\leq a\leq5,\\
       &\hspace{22mm}1\leq e<T,\ \rho_{a,e}=1,\ e\ne r\},\\
 \mathcal R_J^{\rm one}(T)&=\max_{1\leq r\leq N}
                     \mathcal R_{J\setminus\mathcal T^{\ne r}(T)}.
 \end{split}                                                   \tag{MRT.29}
\]
Every branch retains only one weight of the matched family, and
$|\mathcal T^{\ne r}(T)|=2(N-1)$.  For every original retained set $J$,
\[
 \begin{split}
 0\leq\mathcal R_J-\mathcal R_J^{\rm one}(T)
 &\leq2\log7\bigl((N-1)+v_7((N-1)!)\bigr)\\
 &\leq\frac73(N-1)\log7<7\mathcal P(F,T).
 \end{split}                                                   \tag{MRT.30}
\]
If $J$ contains the entire matched family, this difference is also
at least $2(N-1)\log7$.  Thus the retained arithmetic quantity is
strictly smaller for $T\geq3$, by a proved aggregate cost.
An original comparison with coefficient $C'$ now gives
\[
 |\mathcal I_J-\mathcal R_J^{\rm one}(T)|
       \leq E_{\rm old}+(C'+7)\mathcal P(F,T).                  \tag{MRT.31}
\]
The maximum is taken pointwise on one identical host and window.
There is no factor $N$ in the cost, and no separate windows for the
candidate retained weights.  This is a direct replacement of the
old remainder, not the accumulation of the $+6$ and $+7$ alternatives.

\end{theorem}

\begin{proof}
\proofstep{1. Full certification of the fixed class group.}
The field identification from AC-G203 is elementary:
$X^3-269$ is Eisenstein, $[L:\mathbb Q]=6$, and $z=t+\omega$
has the six distinct conjugates $t\omega^i+\omega^a$,
$i=0,1,2$, $a=1,2$.  An equality between two of them would
either identify both indices or put $t$ in $K$.  Its polynomial is
\[
 P(X)=X^6+3X^5+6X^4-531X^3-801X^2+810X+72900.
 \tag{MRT.9}
\]
The retained AC-G203 BNF was previously certified with flag one,
which proved only that the actual class group was a quotient of
$C_7^2$.  A new completed call on an exact copy of that same BNF,
using \texttt{bnfcertify(B,0)}, gives
\[
 B.\mathrm{cyc}=[7,7],\qquad B.\mathrm{no}=49,
                 \qquad \operatorname{bnfcertify}(B,0)=1.
 \tag{MRT.10}
\]
The retained input, complete output, unchanged BNF hash, exit-zero
receipt and actual start/end times bind this execution.  It ran
from 10:58:25.256974 to 10:58:25.520010 UTC on 23 November 2026.
The documented flag-zero contract certifies the full class-group
and unit calculation.  Hence in this proof, unlike the earlier
quotient-only uses, one has the exact group
\[
                         C:=\operatorname{Cl}(L)\simeq C_7^2.
 \tag{MRT.11}
\]
No unfinished degree-eighteen calculation is used.  The complete
applied PARI manual clause and public flag-zero interface are
retained; this is a software certificate, not a new proof of its
underlying algorithm.

For completeness $\operatorname{Cl}(K)=0$ follows from norm
Euclideanity of $\mathbb Z[\omega]$.  Rounding both real coordinates
in the basis $1,\omega$ gives a remainder $a+b\omega$ with
$|a|,|b|\leq1/2$ and norm $a^2-ab+b^2\leq3/4<1$.
The usual least-norm division argument makes every ideal principal.
Let $g(t)=\omega t$ and let $\tau$ fix $t$ and conjugate $\omega$.
Prime-to-seven norm and extension give
$C^{g=1}\simeq\operatorname{Cl}(K)\{7\}=0$.
Indeed $N\circ j=3$ on the base and $j\circ N=1+g+g^2$
on $C$, so their restrictions identify the invariants.
The relation $\tau g\tau^{-1}=g^{-1}$ exchanges the two
nontrivial eigenspaces over $\mathbb F_7$.  Since $C$ has dimension
two, each has dimension one.  Thus for $j=1,2$ the equivariant group
\[
 \{f:C\longrightarrow E/O: f(gx)=\chi^j(g)f(x)\}
                         \simeq O/7O.                         \tag{MRT.12}
\]
This uses the natural class-group action; the displayed equivariance
specifies the orientation and avoids an implicit inverse convention.

\proofstep{2. Principal seven primes and the untwisted strict group.}
Seven splits in $K$, with prime generators $2-\omega$ and
$2-\omega^2$, each of norm seven.  Their ideals are distinct,
since their difference is divisible only above three.
The nonzero cubes modulo seven are $1$ and $-1$, whereas
$269\equiv3\pmod7$.  Thus $X^3-269$ is separable and
irreducible over $\mathbb F_7$.  At each of the two places of $K$
over seven, $L/K$ is unramified of degree three.  There is one
prime above each, with ramification index one; consequently
\[
        \mathfrak P_1=(2-\omega)O_L,
        \qquad \mathfrak P_2=(2-\omega^2)O_L.                  \tag{MRT.13}
\]
Both seven-prime ideal classes are zero.

Put $D_j^0=(E/O)(\chi^j)$, an auxiliary untwisted coefficient.
Restriction from $K$ to $L$ identifies its global first cohomology
with the $C_3$-equivariant continuous homomorphisms from $G_L$
to $E/O$: the positive cohomology of $C_3$ vanishes by averaging.
The same averaging gives injective restriction on every local
decomposition or inertia group.  Therefore the strict condition
away seven is precisely unramifiedness of this homomorphism,
and its strict condition over seven says that it vanishes on the
two unramified Frobenius elements.  The latter is automatic by
(MRT.13).  The equivariant Hilbert class-field Artin map, with
(MRT.12), proves the exact isomorphism
\[
                         \mathcal S_7(K,D_j^0)\simeq O/7O.
 \tag{MRT.14}
\]
This is an auxiliary exponent-zero computation, not the insertion
of a weight zero into the actual source window.

\proofstep{3. Finite coefficients retain every strict local condition.}
Suppose $\rho_{a,e}=1$.  Reduction on $\mu_6$ is injective,
so uniqueness of Teichm\"uller lifts gives the literal identity
$\xi_7^a=\omega_7^e$ on $G_K$.  Thus
\[
             D_{j,a,e}=(E/O)(\chi^j\langle c_7\rangle^{-e}).
 \tag{MRT.15}
\]
In particular all its character values belong to $\mathbb Z_7$.
Its global invariants are zero, because the residual cubic
character is nontrivial.  At each full seven decomposition group,
its residual Frobenius is a primitive cube root.  The difference
from one is an $O$-unit, so its invariants there are zero as well.
Away seven the cyclotomic character is unramified and inertia acts
only by $\chi^j$.  These invariants are either zero or the divisible
module $E/O$.  The identical conclusions hold for $D_j^0$.

For every $n=7^m$, $m\geq1$, apply global and local cohomology to
$0\to D[n]\to D\xrightarrow{n}D\to0$.  The relevant kernel is
$D^{H}/nD^{H}$, which is zero for the global group and for every
local group just considered.  Hence the localization square gives
\[
             \mathcal S_7(K,D[n])\simeq\mathcal S_7(K,D)[n]
                       \qquad(n=7^m,\ m\geq1).                         \tag{MRT.16}
\]
All matched coefficients have $D[7]=D_j^0[7]$.  Equations
(MRT.14)--(MRT.16) prove (MRT.4).  If $7\mid e$, the binomial
identity $(1+7x)^7\equiv1\pmod{49}$ gives
$D[49]=D_j^0[49]$, so $\mathcal S_7(K,D)[49]\simeq O/7O$.
This last group is killed by seven.  Every continuous cocycle into
$E/O$ has finite image and is killed by a power of seven.  An
element of order at least $49$ in the strict group would, after
multiplication by a power of seven, give an element of order exactly
$49$, contradicting the preceding conclusion.  Therefore the
whole strict group is $O/7O$ at every matched weight divisible by
seven.  This step uses the full level-$49$ group, not merely a
bound on the dimension of a socle.

\proofstep{4. The finite local lifting problem over $\mathbb Z_7$.}
To treat all accessible residue classes at once, first take
$O=\mathbb Z_7$.  Let $\Sigma$ consist of the places above
$3,7,269$ and the complex place, and set $G=G_{K,\Sigma}$.
Use full decomposition groups at seven and inertia groups at
the other finite places; denote this finite list of groups by $H_v$.
No cohomological condition occurs at the complex place.
Define the continuous additive character
\[
 \lambda(g)=\frac{\langle c_7(g)\rangle-1}{7}\pmod7.
 \tag{MRT.17}
\]
For $u\in\mathbb F_7$ let $A_{49,j,u}=\mathbb Z/49\mathbb Z$
with action
\[
       g\cdot x=\chi^j(g)(1-7u\lambda(g))x,
       \qquad A_j=\mathbb F_7(\overline\chi^{\,j}).             \tag{MRT.18}
\]
The action is well defined and multiplicative modulo $49$.
For a matched weight $e$, its finite coefficient $D[49]$ is
$A_{49,j,u}$ with $u=e\bmod7$.  We have the exact sequence
\[
 0\longrightarrow A_j\xrightarrow{\iota} A_{49,j,u}
       \xrightarrow{q} A_j\longrightarrow0,
       \qquad \iota(x)=7x,\qquad q(x)=x\pmod7.                \tag{MRT.19}
\]
Global $H^0(G,A_j)$ is zero.  Therefore $\iota_*$ on global
$H^1$ is injective, and multiplication by seven on that group
equals $\iota_*q_*$.  It follows that a strict class for
$A_{49,j,u}$ has order $49$ exactly when its reduction is nonzero
in the one-dimensional strict group for $A_j$.

At every required local group the map
$A_{49,j,u}^{H_v}\to A_j^{H_v}$ is surjective.  At the two full
seven groups both invariants are zero.  At every away inertia
group, either the cubic action is nontrivial and both invariants
are zero, or the entire action is trivial for every $u$ and
surjectivity is immediate.  This observation is needed below
to keep fixed local witnesses without losing any strict lift.

\proofstep{5. An affine obstruction leaves at most one residue.}
Choose a continuous cocycle $z\in Z^1(G,A_j)$ representing a
nonzero strict class, and choose $d_v\in A_j$ with
$z|_{H_v}=d d_v$, where $d$ is the usual inhomogeneous coboundary.
Choose arbitrary continuous lifts $\widetilde z$ and
$\widetilde d_v$ to $\mathbb Z/49\mathbb Z$.
For the action in (MRT.18), divide the global cocycle defect and
the local witness defects by seven and reduce modulo seven:
\[
 B_u=\left(\frac{d_u\widetilde z}{7},
       \left(\frac{\operatorname{res}_v\widetilde z
                      -d_u\widetilde d_v}{7}\right)_v\right)
       \in V:=C^2(G,A_j)\oplus\bigoplus_v C^1(H_v,A_j).
 \tag{MRT.20}
\]
The numerators reduce to zero, so these are well defined.
The action varies affinely in $u$.  Define its slope
$C_{\rm aff}=B_1-B_0$; then $B_u=B_0+uC_{\rm aff}$ in $V$.
Correcting the lifts by $7y$ and $7x_v$ changes their defects by
the fixed linear map
\[
 T:C^1(G,A_j)\oplus\bigoplus_v A_j\longrightarrow V,
 \qquad T(y,(x_v))=(dy,(\operatorname{res}_v y-dx_v)_v).
 \tag{MRT.21}
\]
Consequently $z$ has a strict lift to $A_{49,j,u}$ exactly when
\[
                   [B_0]+u[C_{\rm aff}]=0\quad\hbox{in }V/\operatorname{im}T.
 \tag{MRT.22}
\]
To verify the reverse direction for cohomology classes, first
subtract a lifted global coboundary so that reduction is exactly
$z$.  A local witness may then reduce to $d_v$ plus an invariant.
The local surjectivity proved in Step 4 allows subtracting a lifted
invariant.  The resulting witnesses reduce to the chosen $d_v$,
and therefore are precisely the corrections in (MRT.21).
This also proves that the solvability criterion is independent of
all choices.  A nonzero scalar multiple of $z$ is liftable exactly
when $z$ is liftable.

At $u=0$, (MRT.14) and (MRT.16) show that the finite strict group
is killed by seven.  Step 4 then makes its reduction zero, so
$[B_0]\ne0$.  An affine equation with nonzero constant term in
an $\mathbb F_7$-vector space has at most one solution, and no
solution at zero.  Let $\mathscr E_j$ be this solution set.
For $u\notin\mathscr E_j$, every finite strict class at level
$49$ is killed by seven.  Equations (MRT.4), (MRT.16), and the
torsion argument in Step 3 imply that the entire strict group
is $\mathbb Z/7\mathbb Z$.

Complex conjugation on $K$ sends $\chi$ to $\chi^{-1}$ and
preserves the rational cyclotomic character.  It permutes the
two seven decomposition groups and the remaining inertia list.
It therefore identifies the two strict lifting problems and
gives $\mathscr E_1=\mathscr E_2$.  Define this common set to be
$\mathscr E$.  Changing the faithful cubic character merely
exchanges the two problems.  Changing the faithful $\xi_7$ only
relabels which $a$ is matched to $e$; (MRT.15) is unchanged.
This proves (MRT.3) and the claimed independence.


\proofstep{6. Primitive parents and the same slope at all levels.}
Continue over $\mathbb Z_7$ and write
\[
 D_{j,e}^*=(\mathbb Q_7/\mathbb Z_7)
                    (\chi^j\langle c_7\rangle^{-e}),\qquad
 A_{m,j,e}=D_{j,e}^*[7^m].                                    \tag{MRT.32}
\]
For finite actions an arbitrary integer representative of $e$ is
allowed; for actual groups choose positive representatives.
Binomial expansion shows that $A_{m,j,e}$ depends only on
$e\bmod7^{m-1}$.  Global invariants vanish at every finite level,
and the local invariant reduction maps are all surjective: at the
full seven groups they are zero, and at away inertia they are zero
or the entire, trivially acted-on coefficients.  Equation (MRT.16)
therefore holds at every level, compatibly with inclusions and
multiplication or reduction maps.

An abelian group $M$ killed by $7^m$ with $|M[7]|=7$ is finite
cyclic; no initial finiteness assumption is needed.  Induct on $m$.
The group $(7M)[7]$ is contained in $M[7]$.  If $7M$ is nonzero,
its bounded exponent gives a nonzero element of order seven, so
induction applies and makes $7M$ finite cyclic.  The kernel of
$M\xrightarrow{7}7M$ is $M[7]$, so $|M|=7|7M|\leq7^m$.
The finite abelian group structure theorem, or a lift of a generator
of $7M$, makes $M$ cyclic.  The case $7M=0$ is immediate.
Thus every finite-coefficient strict group
$\mathcal S_7(K,A_{m,j,e})$ is cyclic of order at most $7^m$.

Global $H^0(G,A_{k,j,e})=0$ for all $k$ makes the inclusions on
$H^1$ in every sequence
$0\to A_{r,j,e}\to A_{r+s,j,e}\to A_{s,j,e}\to0$ injective.
Multiplication by seven on $H^1(A_{m,j,e})$ is reduction to
$H^1(A_{m-1,j,e})$ followed by this injective inclusion.
Consequently an order-$7^m$ class reduces to order $7^{m-1}$,
and a lift of an order-$7^{m-1}$ class has order $7^m$.
In particular a primitive class reduces nontrivially to $A_j$.

Fix the original residual cocycle $z$ and local witnesses $d_v$
from Step 5.  Whenever a parent strict group at level $m$ contains
an order-$7^m$ class, it is cyclic of that full order.
Scale a generator by a unit so that its residual cohomology class
is $[z]$.  Subtracting a lifted global coboundary makes its cocycle
reduce exactly to $z$.  Its local witnesses can be made to reduce
exactly to $d_v$ by the surjectivity of invariant reduction.
Every other primitive parent differs by a unit, so normalization
has discarded no possible primitive lift.

If $\mathscr E$ is empty, no level-$49$ primitive class exists,
and no higher primitive class can exist.  Suppose instead that
$\mathscr E$ is nonempty.  In (MRT.22), $[B_0]\ne0$, so the slope
$[C_{\rm aff}]$ is nonzero.  For $n\geq3$, suppose a parent residue
$e_0\bmod7^{n-2}$ admits a primitive level-$7^{n-1}$ class.
Normalize a parent cocycle and its witnesses as above and lift them
arbitrarily to $\mathbb Z/7^n$.  The seven children have exponents
$e=e_0+u7^{n-2}$, $u\in\mathbb F_7$.  Their scalar actions satisfy
\[
 \langle c_7(g)\rangle^{-u7^{n-2}}
              \equiv1-u7^{n-1}\lambda(g)\pmod{7^n}.
 \tag{MRT.33}
\]
Divide their global cocycle and local witness defects by $7^{n-1}$
and reduce modulo seven.  Correcting the lifts by $7^{n-1}y$ and
$7^{n-1}x_v$ again changes defects by the \emph{same} map $T$
in (MRT.21).  Moreover the slope cochains are exactly
\[
 \left(-\overline{\chi}^{\,j}(g)\lambda(g)z(h),\quad
           (\overline{\chi}^{\,j}(h)\lambda(h)d_v)_v\right).
 \tag{MRT.34}
\]
These are the slope cochains of (MRT.20), because the normalized
parent cocycle, witnesses and action reduce to the same $z,d_v$
and $\overline{\chi}^{\,j}$.  The slope class is therefore precisely
$[C_{\rm aff}]\ne0$, not a new, possibly vanishing class.
Only the constant defect may change with the parent.

The child lifting criterion is the new constant plus
$u[C_{\rm aff}]=0$ in the same quotient $V/\operatorname{im}T$.
For its necessity, scale a primitive child's parent reduction to
the fixed generator, lift the required global coboundary, and
adjust each witness using
$A_{n,j,e}^{H_v}\twoheadrightarrow A_{n-1,j,e}^{H_v}$.
For sufficiency, a solution of the defect equation supplies the
cocycle and every local witness.  Thus at most one child survives.
The initial level already has at most one parent by Step 5.
Induction proves that
\[
 \begin{split}
 E_n=\{r\in\mathbb Z/7^{n-1}\mathbb Z:\quad&
       \mathcal S_7(K,A_{n,j,r})\\
       &\text{ contains an element of order }7^n\}
 \end{split}
 \tag{MRT.35}
\]
is empty or a singleton.  Every higher residue reduces to its
parent, and an empty level has no descendants.  Conjugation exchanges
$j=1,2$ at all finite levels and preserves the rational cyclotomic
action and local list.  Hence the entire nested sequence is common
to both cubic characters.  No Iwasawa presentation or characteristic
power series has been assumed.

\proofstep{7. Original coefficient rings and the actual host group.}
The preceding actions are all defined over $\mathbb Z_7$.
For any original $O$, finite freeness over $\mathbb Z_7$ identifies
the continuous cochain complexes and their localization maps
with the base complexes tensored with $O$.  Flatness, or an
explicit $\mathbb Z_7$-basis of $O$, identifies their strict
kernels as well.  In particular the bound holds on the actual
ramified ring, with $O/49O$ at level $49$, rather than a replacement
by a squared uniformizer quotient.  This proves (MRT.4)--(MRT.5)
for every original $O$, and $\#(O/7O)=7^b$ proves the normalized
cost exactly.

It remains to justify equality, not merely injection, for (FC.2).
Let $\Sigma_7^{\rm host}$ be its actual localization set, which
contains the places over $3,7,269$ and every ramified place of $D$.
Outside this set inertia acts trivially on $D$.  A global strict
cocycle restricts there to the zero cohomology class, hence to the
zero homomorphism: with trivial action every local coboundary is
zero.  The cocycle conjugation formula preserves this literal zero.
The closed normal kernel of $G_K\to G_{K,\Sigma_7^{\rm host}}$
is generated by these inertia groups and acts trivially on $D$.
The cocycle restricts to a continuous homomorphism on that kernel,
vanishes on its normal generators, and therefore vanishes everywhere
on it.  It descends uniquely.  Conversely inflation preserves every
full-seven and away-inertia kernel.  This proves equality of the
global strict group and the actual (FC.2) group, with all extra host
inertia conditions retained.  The same argument justified the
minimal set $\Sigma$ used in Steps 4--5.

The retained finiteness statement (FC.2), equivalently the explicit
finite-sum scope in (CP.3), applies at every positive actual weight.
Use any one actual coefficient ring as a finite free, faithfully
flat extension of $\mathbb Z_7$.  The cochain comparison just proved
then implies that $\mathcal S_7(K,D_{j,e}^*)$ is finite as well;
this does not silently assume that the actual ring was $\mathbb Z_7$.
Its seven-socle has dimension one, so it is cyclic of order
$7^{s(e)}$ for a positive integer $s(e)$.  Conjugation makes this
integer independent of $j$.  The exact finite-level comparisons give
\[
 s(e)\geq n\quad\Longleftrightarrow\quad
                   e\bmod7^{n-1}\in E_n\qquad(n\geq2).
 \tag{MRT.36}
\]
Finite free base change now gives (MRT.27) for every original $O$.
If the nested sequence never terminates, it determines a single
$\alpha\in\mathbb Z_7$ with
$s(e)=1+v_7(e-\alpha)$ for positive integers $e$.
This $\alpha$ cannot be a positive integer, by the just established
finiteness.  If the first empty set is $E_{m+1}$ with $m\geq2$, then for any lift
$\alpha$ of the last residue $E_m$ one has
$s(e)=\min(m,1+v_7(e-\alpha))$; for $m=1$ the formula simply says
$s(e)=1$ and any $\alpha$ may be chosen.  Here $v_7(0)=+\infty$.
These descriptions identify no numerical $\alpha$ or stopping level,
and no quantitative approximation bound for $\alpha$ is used.


\proofstep{8. The complete family, its cost, and the remaining gate.}
Reduction of the actual faithful $\xi_7$ is a generator of the
six rational characters.  For each positive $e$ exactly one $a$
has $\rho_{a,e}=1$ in its original coefficient embedding.  There
are two cubic powers, so deleting precisely the exceptional
residue weights gives (MRT.6).  In particular the explicitly known
subfamily $7\mid e$ has $2\lfloor N/7\rfloor$ rows.  It is
included in this single larger theorem, not counted as a separate
intermediate reduction.

The coefficient-seven families of AC-G213 have residual order
three.  Those of AC-G217 have residual order two when the two-power
factor is trivial, and residual order one only with a nontrivial
two-power quadratic factor.  AC-G211 also has such a nontrivial
quadratic factor.  The growing-prime families AC-G203, AC-G207,
AC-G209 and AC-G215 exclude seven.  Thus (MRT.6) lies in the
previously retained pure residual-matched cubic part, and its
positive terms have not been removed by those results.

Every term in (MRT.6) is a mixed term of (CP.3), with exact value
$\log7$ by (MRT.5).  Summing over an arbitrary $J$ proves the
equality in (MRT.7), without a stability assumption on $J$.
Since $D,\Delta,R\geq1$ in (CP.2),
\[
 \mathcal P(F,T)\geq T\log2,\qquad
 \Gamma_J\leq2(T-1)\log7<6T\log2\leq6\mathcal P(F,T),
 \tag{MRT.23}
\]
where $49<64$ gives the strict middle inequality.
The triangle inequality now proves (MRT.8).
For an arithmetic threshold $B$ the exact equivalent remainder
condition is $\mathcal R_{J'}\leq B-\Gamma_J$; a uniform sufficient
condition is $\mathcal R_{J'}\leq B-6\mathcal P$.
This is a controlled threshold adjustment, not an assertion that
the previous slack is numerically unchanged.  The existing full
finite-gluing, torsion and critical comparison errors are counted
once, and the same prescribed window and original triples are used.
For the explicit finite-choice form, nonnegativity gives
$\mathcal R_{J\setminus\mathcal A_r}\leq\mathcal R_J$ for every $r$.
If $\mathscr E=\{r_*\}$, choose that $r_*$; if $\mathscr E$ is
empty, choose any $r_*\in\{1,\ldots,6\}$.
All deleted rows in $\mathcal A_{r_*}$ then have value $\log7$.
Hence
\[
 \begin{split}
 0\leq\mathcal R_J-\mathcal R^{\max}_J
 &\leq\mathcal R_J-\mathcal R_{J\setminus\mathcal A_{r_*}}\\
 &=|J\cap\mathcal A_{r_*}|\log7\leq2(T-1)\log7.
 \end{split}                                                   \tag{MRT.26}
\]
This proves (MRT.25) and its comparison consequence without computing
or declaring a value for the exceptional residue.  The maximum is
taken pointwise on one identical host and window; no factor six
multiplies the finite-gluing error or the known cost.

To obtain the stronger single-weight bound, choose $e_*$ with maximal
$s(e)$ on $1\leq e\leq N$.  For each nonempty level set
$\{e:s(e)\geq n\}$ this weight is a member.  For $n\geq2$,
(MRT.35)--(MRT.36) put that set in one residue class modulo
$7^{n-1}$, containing at most $1+\lfloor(N-1)/7^{n-1}\rfloor$
integers in the interval.  After removing $e_*$, its size is at
most $\lfloor(N-1)/7^{n-1}\rfloor$.  At level one it is exactly
$N-1$.  Summing the finite layer counts proves (MRT.28).
The factorial identity counts multiples of each $7^k$; it also
holds at $N=1$, with $0!=1$.  The geometric-series bound is
$\sum_{k\geq0}(N-1)/7^k=7(N-1)/6$.

Every matched weight has two actual cubic labels, with the same
normalized size $s(e)\log7$ on every original $O$.
Nonnegativity gives
$\mathcal R_{J\setminus\mathcal T^{\ne r}}\leq\mathcal R_J$
for each $r$.  Choosing $r=e_*$ shows
\[
 \begin{split}
 0\leq\mathcal R_J-\mathcal R_J^{\rm one}(T)
 &\leq\mathcal R_J-
           \mathcal R_{J\setminus\mathcal T^{\ne e_*}(T)}\\
 &\leq2\log7\sum_{\substack{1\leq e\leq N\\e\ne e_*}}s(e).
 \end{split}                                                   \tag{MRT.37}
\]
This remains true if $J$ omits either or both labels at some weights.
Together with (MRT.28) and
\[
 \frac73(N-1)\log7<7T\log2\leq7\mathcal P(F,T),
 \tag{MRT.38}
\]
it proves (MRT.30).  The strict inequality follows from $7<8$.
If all matched rows belong to $J$, each candidate removes exactly
$2(N-1)$ rows of normalized size at least $\log7$, so the minimum
removed cost, which equals $\mathcal R_J-\mathcal R_J^{\rm one}$,
is at least $2(N-1)\log7$.  This is a literal change of the retained
original arithmetic quantity, not a renaming of a hypothesis.

The triangle inequality proves (MRT.31).  In particular
$\mathcal R_J^{\rm one}\leq B-7\mathcal P$ is a uniform sufficient
condition for the old arithmetic threshold $\mathcal R_J\leq B$.
The single-weight maximum replaces the old remainder directly;
one does not first charge the six-residue cost and then charge this
cost again.  Only one full finite-gluing and critical comparison is
used on the identical $W$, and no maximum requires separately
existential windows.  The sole matched weight retained by a maximizing
branch is unbounded by this theorem; all other old arithmetic
families remain exactly as before.

For $T=\lfloor u^{399}\rfloor$, (CP.5) puts the added cost inside
the existing polynomial envelope and hence below its eventual
exponential slack.  For the single-weight formulation, one maximal
matched weight and
every other retained family remain in the arithmetic remainder.
\end{proof}

\paragraph{Finite-\'etale and source scope.}
The finite modules $D[7^m]$ are the original finite Artin--Tate
sheaves on $\operatorname{Spec}O_K[1/(3\cdot7\cdot269)]$.
The residual cubic cover is $L/K$, and the full completion and
inertia restriction kernels are exactly the strict local conditions.
The finite-\'etale dictionary is the retained SGA~1, Expos\'e~V,
sections 4 and 7 scope already used in the fixed-base audit; no new
archive reading or motive realization is asserted.
The deep Hilbert class-field identification is used as the explicit
retained primary theorem.  Its norm and conjugation applications
are checked in the complete applied Sharifi proofs.
The fixed-degree and $\mu=0$ obstruction AC-K28 remains valid:
the present bound uses an actual full class certificate and the
actual finite-level deformations and their repeated nonzero slope,
not generic Iwasawa degree or a regulator hypothesis.
In particular it does not bound the one
retained closest weight by such an assumption.

\paragraph{Contribution and completion scope.}
The parent proposed the principal-prime obstruction test, the
actual-host descent check and the finite maximum formulation.
The author completed the full flag-zero
certificate and the finite-$49$ transfer.  The independent reviewer
proposed the affine and repeated-lifting enlargements; their complete
cochain and local-witness proofs are given above.  The first reviewed
finite-$49$ version remains immutable in its original directory; this
comprehensive revision is separately frozen with actual chronology.
It is the selected version of the same finding, not another helper or
a separately banked narrow result.  Earlier quotient-only certificates,
unsuccessful stronger
routes, and the unfinished degree-eighteen run retain their original
scopes.  No class-socle diagnostic, finite calculation, or packaging
step is counted as a separate arithmetic-gate reduction.
